% !TeX root = part3.tex

\documentclass[../final.tex]{subfiles}

\begin{document}


% ============================================================
% Семинар 3
% ============================================================

\practice{3}{19.09.2026}{Колебания решётки}


% ============================================================
% Задача 1
% ============================================================

\task[1]{Линейная цепочка с двумя атомами в элементарной ячейке}

\condition

Рассмотреть бесконечную одномерную цепочку,
состоящую из чередующихся атомов масс $m$ и $M$,
соединённых одинаковыми пружинами жёсткости $\alpha$.

Расстояние между соседними атомами равно $a$,
поэтому период элементарной ячейки, содержащей два атома,
равен $2a$.

Обозначим смещения атомов в $n$-й элементарной ячейке через
%
\begin{equation*}
    u_n(t),
    \qquad
    v_n(t).
\end{equation*}
%
Необходимо определить собственные колебания цепочки,
получить закон дисперсии $\omega(k)$,
выделить акустическую и оптическую ветви
и исследовать их поведение в центре и на границе
первой зоны Бриллюэна.


% ------------------------------------------------------------
% РИСУНОК 1
% ------------------------------------------------------------

\begin{rbox}[left=18pt,right=18pt,top=14pt,bottom=10pt]{[]}
    \centering

    \begin{tikzpicture}[
        >=Latex,
        every node/.style={text=rtext}
    ]

        % Продолжение цепочки
        \node at (-6.0,0) {$\ldots$};
        \node at ( 6.0,0) {$\ldots$};

        % Пружины
        \foreach \x in {-5,-3,-1,1,3}{
            \draw[
                draw=rtext,
                line width=0.8pt
            ]
            plot[
                domain=0:1,
                samples=70,
                variable=\t
            ]
            (
                {\x+0.30+1.40*\t},
                {0.11*sin(1440*\t)}
            );
        }

        % Атомы m
        \foreach \x in {-5,-1,3}{
            \filldraw[
                fill=rc3,
                draw=rtext,
                line width=0.9pt
            ]
            (\x,0) circle (0.28);

            \node[above=5pt] at (\x,0) {$m$};
        }

        % Атомы M
        \foreach \x in {-3,1,5}{
            \filldraw[
                fill=rbgsecondary,
                draw=rtext,
                line width=0.9pt
            ]
            (\x,0) circle (0.28);

            \node[above=5pt] at (\x,0) {$M$};
        }

        % Индексы центральных ячеек
        \node[below=5pt] at (-5,0) {$u_{n-1}$};
        \node[below=5pt] at (-3,0) {$v_{n-1}$};

        \node[below=5pt] at (-1,0) {$u_n$};
        \node[below=5pt] at ( 1,0) {$v_n$};

        \node[below=5pt] at ( 3,0) {$u_{n+1}$};
        \node[below=5pt] at ( 5,0) {$v_{n+1}$};

        % Жёсткости
        \foreach \x in {-4,-2,0,2,4}{
            \node[above] at (\x,0.18) {$\alpha$};
        }

        % Период элементарной ячейки
        \draw[
            <->,
            draw=rtext,
            line width=0.8pt
        ]
        (-1,-1.0)
        --
        (3,-1.0)
        node[midway,below] {$2a$};

    \end{tikzpicture}

    \captionsetup{
        font=footnotesize,
        labelfont=bf,
        skip=5pt
    }

    \captionof{figure}{
        Двухатомная одномерная цепочка.
        Атомы масс $m$ и $M$ чередуются,
        соседние атомы соединены пружинами жёсткости $\alpha$.
        Период элементарной ячейки равен $2a$.
    }

    \label{fig:diatomic_chain}
\end{rbox}


\solution

\subsection*{Уравнения движения}

Рассмотрим атом массы $m$,
смещение которого в $n$-й элементарной ячейке равно $u_n$.

С правой стороны он связан с атомом,
имеющим смещение $v_n$,
а с левой стороны --- с атомом $v_{n-1}$.

Поэтому
%
\begin{equation*}
    m\ddot u_n
    =
    \alpha(v_n-u_n)
    -
    \alpha(u_n-v_{n-1}).
\end{equation*}
%

Для атома массы $M$ со смещением $v_n$
аналогично получаем
%
\begin{equation*}
    M\ddot v_n
    =
    \alpha(u_{n+1}-v_n)
    -
    \alpha(v_n-u_n).
\end{equation*}
%

Таким образом,
%
\begin{align*}
    m\ddot u_n
    &=
    \alpha
    \left(
        v_n+v_{n-1}-2u_n
    \right),
    \\
    M\ddot v_n
    &=
    \alpha
    \left(
        u_{n+1}+u_n-2v_n
    \right).
\end{align*}


% ------------------------------------------------------------
% Нормальная мода
% ------------------------------------------------------------

\subsection*{Нормальная мода}

Поскольку период элементарной ячейки равен $2a$,
ищем решение в виде
%
\begin{equation*}
    u_n
    =
    U_0
    e^{-i\omega t+ik\,2an},
\end{equation*}
%
\begin{equation*}
    v_n
    =
    V_0
    e^{-i\omega t+ik\,2an}.
\end{equation*}
%

Для соседних ячеек
%
\begin{equation*}
    v_{n-1}
    =
    v_n e^{-2ika},
    \qquad
    u_{n+1}
    =
    u_n e^{2ika}.
\end{equation*}
%

Кроме того,
%
\begin{equation*}
    \ddot u_n=-\omega^2u_n,
    \qquad
    \ddot v_n=-\omega^2v_n.
\end{equation*}


% ------------------------------------------------------------
% Первое уравнение
% ------------------------------------------------------------

\subsection*{Уравнение для амплитуды $U_0$}

Подставим нормальную моду в первое уравнение движения:
%
\begin{equation*}
    -m\omega^2
    U_0e^{-i\omega t+ik\,2an}
    =
    \alpha
    V_0e^{-i\omega t+ik\,2an}
    -
    \alpha
    U_0e^{-i\omega t+ik\,2an}
\end{equation*}
%
\begin{equation*}
    \hspace{3cm}
    -
    \alpha
    \left(
        U_0e^{-i\omega t+ik\,2an}
        -
        V_0e^{-i\omega t+ik\,2a(n-1)}
    \right).
\end{equation*}
%

Сокращая общий множитель
%
\begin{equation*}
    e^{-i\omega t+ik\,2an},
\end{equation*}
%
получаем
%
\begin{equation*}
    -m\omega^2U_0
    =
    \alpha V_0
    -
    \alpha U_0
    -
    \alpha
    \left(
        U_0-V_0e^{-2ika}
    \right).
\end{equation*}
%

Следовательно,
%
\begin{equation*}
    \left(
        m\omega^2-2\alpha
    \right)U_0
    +
    \alpha
    \left(
        1+e^{-2ika}
    \right)V_0
    =
    0.
\end{equation*}


% ------------------------------------------------------------
% Второе уравнение
% ------------------------------------------------------------

\subsection*{Уравнение для амплитуды $V_0$}

Для второго атома
%
\begin{equation*}
    -M\omega^2V_0
    =
    \alpha
    \left(
        U_0e^{2ika}-V_0
    \right)
    -
    \alpha
    \left(
        V_0-U_0
    \right).
\end{equation*}
%

Отсюда
%
\begin{equation*}
    \alpha
    \left(
        1+e^{2ika}
    \right)U_0
    +
    \left(
        M\omega^2-2\alpha
    \right)V_0
    =
    0.
\end{equation*}


% ------------------------------------------------------------
% Матричная форма
% ------------------------------------------------------------

\subsection*{Матричная форма}

Полученную систему можно записать в виде
%
\begin{equation*}
    \begin{pmatrix}
        m\omega^2-2\alpha
        &
        \alpha(1+e^{-2ika})
        \\[2mm]
        \alpha(1+e^{2ika})
        &
        M\omega^2-2\alpha
    \end{pmatrix}
    \begin{pmatrix}
        U_0
        \\
        V_0
    \end{pmatrix}
    =
    0.
\end{equation*}
%

Нетривиальное решение существует,
если определитель матрицы равен нулю:
%
\begin{equation*}
    \left(
        m\omega^2-2\alpha
    \right)
    \left(
        M\omega^2-2\alpha
    \right)
    -
    \alpha^2
    \left(
        1+e^{-2ika}
    \right)
    \left(
        1+e^{2ika}
    \right)
    =
    0.
\end{equation*}


% ------------------------------------------------------------
% Дисперсионное уравнение
% ------------------------------------------------------------

\subsection*{Дисперсионное уравнение}

Раскрываем первое произведение:
%
\begin{equation*}
    mM\omega^4
    -
    2\alpha(m+M)\omega^2
    +
    4\alpha^2.
\end{equation*}
%

Для второго произведения имеем
%
\begin{align*}
    \left(
        1+e^{-2ika}
    \right)
    \left(
        1+e^{2ika}
    \right)
    &=
    2
    +
    e^{2ika}
    +
    e^{-2ika}
    \\
    &=
    2
    +
    2\cos 2ka.
\end{align*}
%

Поэтому условие на определитель принимает вид
%
\begin{equation*}
    mM\omega^4
    -
    2\alpha(m+M)\omega^2
    +
    4\alpha^2
    -
    2\alpha^2
    \left(
        1+\cos 2ka
    \right)
    =
    0.
\end{equation*}
%

Используем
%
\begin{equation*}
    1+\cos2ka
    =
    2\cos^2ka.
\end{equation*}
%

Тогда
%
\begin{equation*}
    4\alpha^2
    -
    4\alpha^2\cos^2ka
    =
    4\alpha^2\sin^2ka.
\end{equation*}
%

Окончательно получаем
%
\begin{equation*}
    mM\omega^4
    -
    2\alpha(m+M)\omega^2
    +
    4\alpha^2\sin^2ka
    =
    0.
\end{equation*}


% ------------------------------------------------------------
% Две ветви
% ------------------------------------------------------------

\subsection*{Акустическая и оптическая ветви}

Полученное уравнение является квадратным
относительно $\omega^2$.

Отсюда
%
\begin{equation*}
    \omega_{\pm}^{\,2}(k)
    =
    \frac{
        \alpha(m+M)
        \pm
        \sqrt{
            \alpha^2(m+M)^2
            -
            4mM\alpha^2\sin^2ka
        }
    }{
        mM
    }.
\end{equation*}
%

Или, вынося общий множитель,
%
\begin{equation*}
    \omega_{\pm}^{\,2}(k)
    =
    \frac{\alpha(m+M)}{mM}
    \left[
        1
        \pm
        \sqrt{
            1
            -
            \frac{4mM}{(m+M)^2}
            \sin^2ka
        }
    \right].
\end{equation*}
%

Ветвь со знаком минус является
\rconcept{акустической ветвью},
%
\begin{equation*}
    \omega_-^{\,2}(k)
    =
    \frac{\alpha(m+M)}{mM}
    \left[
        1
        -
        \sqrt{
            1
            -
            \frac{4mM}{(m+M)^2}
            \sin^2ka
        }
    \right],
\end{equation*}
%
а ветвь со знаком плюс ---
\rconcept{оптической ветвью},
%
\begin{equation*}
    \omega_+^{\,2}(k)
    =
    \frac{\alpha(m+M)}{mM}
    \left[
        1
        +
        \sqrt{
            1
            -
            \frac{4mM}{(m+M)^2}
            \sin^2ka
        }
    \right].
\end{equation*}


% ------------------------------------------------------------
% Одинаковые массы
% ------------------------------------------------------------

\subsection*{Предел одинаковых масс}

Если
%
\begin{equation*}
    m=M,
\end{equation*}
%
то
%
\begin{equation*}
    \frac{4mM}{(m+M)^2}=1.
\end{equation*}
%

Поэтому
%
\begin{equation*}
    \omega_\pm^{\,2}
    =
    \frac{2\alpha}{m}
    \left(
        1
        \pm
        \sqrt{
            1-\sin^2ka
        }
    \right).
\end{equation*}
%

В первой зоне Бриллюэна
%
\begin{equation*}
    -\frac{\pi}{2a}
    \leq
    k
    \leq
    \frac{\pi}{2a},
\end{equation*}
%
поэтому $\cos ka\geq0$ и
%
\begin{equation*}
    \sqrt{1-\sin^2ka}
    =
    \cos ka.
\end{equation*}
%

Следовательно,
%
\begin{equation*}
    \omega_\pm^{\,2}
    =
    \frac{2\alpha}{m}
    \left(
        1\pm\cos ka
    \right).
\end{equation*}


% ------------------------------------------------------------
% k -> 0
% ------------------------------------------------------------

\subsection*{Длинноволновый предел}

Рассмотрим
%
\begin{equation*}
    k\to0.
\end{equation*}
%

Для оптической ветви
%
\begin{equation*}
    \sin ka\to0,
\end{equation*}
%
поэтому
%
\begin{equation*}
    \omega_+^{\,2}(0)
    =
    \frac{2\alpha(m+M)}{mM}.
\end{equation*}
%

Для акустической ветви
%
\begin{equation*}
    \omega_-^{\,2}(0)=0.
\end{equation*}
%

Чтобы получить её поведение при малых $k$,
используем разложение
%
\begin{equation*}
    \sqrt{1-x}
    \simeq
    1-\frac{x}{2},
    \qquad
    x\ll1.
\end{equation*}
%

Здесь
%
\begin{equation*}
    x
    =
    \frac{4mM}{(m+M)^2}
    \sin^2ka.
\end{equation*}
%

Тогда
%
\begin{align*}
    \omega_-^{\,2}
    &\simeq
    \frac{\alpha(m+M)}{mM}
    \frac{2mM}{(m+M)^2}
    \sin^2ka
    \\
    &=
    \frac{2\alpha}{m+M}
    \sin^2ka.
\end{align*}
%

При $ka\ll1$
%
\begin{equation*}
    \sin ka\simeq ka,
\end{equation*}
%
и поэтому
%
\begin{equation*}
    \omega_-^{\,2}
    \simeq
    \frac{2\alpha a^2}{m+M}
    k^2.
\end{equation*}
%

То есть
%
\begin{equation*}
    \omega_-(k)\propto |k|.
\end{equation*}
%

Именно линейное поведение при малых $k$
характерно для акустической ветви.


% ------------------------------------------------------------
% Граница зоны
% ------------------------------------------------------------

\subsection*{Граница первой зоны Бриллюэна}

Поскольку период элементарной ячейки равен $2a$,
границы первой зоны Бриллюэна имеют вид
%
\begin{equation*}
    k
    =
    \pm\frac{\pi}{2a}.
\end{equation*}
%

При
%
\begin{equation*}
    k=\frac{\pi}{2a}
\end{equation*}
%
имеем
%
\begin{equation*}
    \sin^2ka=1.
\end{equation*}
%

Поэтому
%
\begin{equation*}
    \omega_\pm^{\,2}
    =
    \frac{\alpha(m+M)}{mM}
    \left[
        1
        \pm
        \sqrt{
            1-\frac{4mM}{(m+M)^2}
        }
    \right].
\end{equation*}
%

Под корнем
%
\begin{equation*}
    1-\frac{4mM}{(m+M)^2}
    =
    \frac{(M-m)^2}{(M+m)^2}.
\end{equation*}
%

Далее, как на схеме в конспекте,
считаем
\rassumption{$M>m$}.
Тогда
%
\begin{equation*}
    \sqrt{
        1-\frac{4mM}{(m+M)^2}
    }
    =
    \frac{M-m}{M+m}.
\end{equation*}
%

Для верхней ветви
%
\begin{align*}
    \omega_+^{\,2}
    &=
    \frac{\alpha(m+M)}{mM}
    \left(
        1+\frac{M-m}{M+m}
    \right)
    \\
    &=
    \frac{\alpha(m+M)}{mM}
    \frac{2M}{M+m}
    \\
    &=
    \frac{2\alpha}{m}.
\end{align*}
%

Для нижней ветви
%
\begin{align*}
    \omega_-^{\,2}
    &=
    \frac{\alpha(m+M)}{mM}
    \left(
        1-\frac{M-m}{M+m}
    \right)
    \\
    &=
    \frac{\alpha(m+M)}{mM}
    \frac{2m}{M+m}
    \\
    &=
    \frac{2\alpha}{M}.
\end{align*}
%

Таким образом, на границе зоны
%
\begin{equation*}
    \omega_-
    =
    \sqrt{\frac{2\alpha}{M}},
    \qquad
    \omega_+
    =
    \sqrt{\frac{2\alpha}{m}}.
\end{equation*}


% ------------------------------------------------------------
% РИСУНОК 2
% ------------------------------------------------------------

\begin{rbox}[left=18pt,right=18pt,top=14pt,bottom=12pt]{[]}
    \centering

    \begin{tikzpicture}[
        x=1.1cm,
        y=1.0cm,
        >=Latex,
        every node/.style={text=rtext}
    ]

        % Оси
        \draw[
            ->,
            draw=rtext,
            line width=0.9pt
        ]
        (-4.2,0)
        --
        (4.4,0)
        node[right] {$k$};

        \draw[
            ->,
            draw=rtext,
            line width=0.9pt
        ]
        (0,-0.15)
        --
        (0,5.2)
        node[above] {$\omega$};

        % Границы зоны
        \draw[
            dashed,
            draw=rtext,
            opacity=0.55
        ]
        (-3.5,0)
        --
        (-3.5,4.7);

        \draw[
            dashed,
            draw=rtext,
            opacity=0.55
        ]
        (3.5,0)
        --
        (3.5,4.7);

        \node[below=5pt] at (-3.5,0)
            {$-\dfrac{\pi}{2a}$};

        \node[below=5pt] at (3.5,0)
            {$\dfrac{\pi}{2a}$};

        % Запрещённая область
        \fill[
            red3,
            opacity=0.10
        ]
        (-3.5,2.15)
        rectangle
        (3.5,3.15);

        \node[
            text=red3
        ]
        at (0,2.65)
        {\small запрещённая зона};

        % Акустическая ветвь
        \draw[
            draw=rc3,
            line width=1.3pt
        ]
        (-3.5,2.15)
        .. controls (-2.1,1.05) and (-0.9,0.20) ..
        (0,0)
        .. controls (0.9,0.20) and (2.1,1.05) ..
        (3.5,2.15);

        % Оптическая ветвь
        \draw[
            draw=p3,
            line width=1.3pt
        ]
        (-3.5,3.15)
        .. controls (-2.2,3.65) and (-1.0,4.15) ..
        (0,4.30)
        .. controls (1.0,4.15) and (2.2,3.65) ..
        (3.5,3.15);

        % Подписи граничных частот
        \draw[
            dashed,
            draw=rc3,
            opacity=0.7
        ]
        (0,2.15)
        --
        (4.0,2.15);

        \draw[
            dashed,
            draw=p3,
            opacity=0.7
        ]
        (0,3.15)
        --
        (4.0,3.15);

        \node[
            right,
            text=rc3
        ]
        at (3.55,2.15)
        {$\sqrt{\dfrac{2\alpha}{M}}$};

        \node[
            right,
            text=p3
        ]
        at (3.55,3.15)
        {$\sqrt{\dfrac{2\alpha}{m}}$};

        % Подписи ветвей
        \node[
            text=rc3
        ]
        at (-1.85,0.65)
        {\small акустическая};

        \node[
            text=p3
        ]
        at (-1.6,4.05)
        {\small оптическая};

    \end{tikzpicture}

    \captionsetup{
        font=footnotesize,
        labelfont=bf,
        skip=5pt
    }

    \captionof{figure}{
        Схематический фононный спектр двухатомной цепочки
        при $M>m$.
        Между акустической и оптической ветвями
        возникает запрещённый интервал частот.
    }

    \label{fig:diatomic_dispersion}
\end{rbox}


Между граничными частотами
%
\begin{equation*}
    \sqrt{\frac{2\alpha}{M}}
    <
    \omega
    <
    \sqrt{\frac{2\alpha}{m}}
\end{equation*}
%
нет нормальных колебаний идеальной цепочки.

Этот интервал представляет собой
\rconcept{запрещённую зону фононного спектра}.


\answer

Дисперсионное соотношение двухатомной цепочки имеет вид
%
\begin{equation*}
    \omega_{\pm}^{\,2}(k)
    =
    \frac{\alpha(m+M)}{mM}
    \left[
        1
        \pm
        \sqrt{
            1
            -
            \frac{4mM}{(m+M)^2}
            \sin^2ka
        }
    \right].
\end{equation*}
%

Нижняя ветвь является акустической,
верхняя --- оптической.

При $M>m$ на границе первой зоны Бриллюэна
%
\begin{equation*}
    \omega_-
    =
    \sqrt{\frac{2\alpha}{M}},
    \qquad
    \omega_+
    =
    \sqrt{\frac{2\alpha}{m}},
\end{equation*}
%
а между этими частотами находится запрещённая зона.


% ============================================================
% Задача 2
% ============================================================

\task[2]{Дефект решётки: пружина другой жёсткости}

\condition

Рассмотреть двухатомную одномерную цепочку
из предыдущей задачи.

Все соседние атомы соединены пружинами жёсткости $\alpha$,
кроме одной связи между атомами $u_0$ и $v_0$.
Жёсткость этой дефектной пружины равна
%
\begin{equation*}
    \alpha_0\neq\alpha.
\end{equation*}
%

Исследовать возможность существования
локализованного колебательного состояния.

Показать, что частота связанного состояния,
если такое состояние существует,
находится в запрещённой зоне фононного спектра.


% ------------------------------------------------------------
% РИСУНОК 3
% ------------------------------------------------------------

\begin{rbox}[left=18pt,right=18pt,top=14pt,bottom=10pt]{[]}
    \centering

    \begin{tikzpicture}[
        >=Latex,
        every node/.style={text=rtext}
    ]

        \node at (-6.0,0) {$\ldots$};
        \node at ( 6.0,0) {$\ldots$};

        % Четыре обычные связи; концы пружин лежат на границах атомов.
        \foreach \x in {-5,-3,1,3}{
            \draw[
                draw=rtext,
                line width=0.8pt
            ]
            ({\x+0.28},0)--({\x+0.40},0)
            --
            plot[
                domain=0:1,
                samples=70,
                variable=\t
            ]
            (
                {\x+0.40+1.20*\t},
                {0.11*sin(1440*\t)}
            )
            --({\x+1.72},0);

            \node at ({\x+1},0.60) {$\alpha$};
        }

        % Единственная дефектная связь: от u_0 при x=-1 до v_0 при x=1.
        \draw[
            draw=red4,
            line width=1.2pt
        ]
        (-0.72,0)--(-0.60,0)
        --
        plot[
            domain=0:1,
            samples=70,
            variable=\t
        ]
        (
            {-0.60+1.20*\t},
            {0.11*sin(1440*\t)}
        )
        --(0.72,0);

        \node[text=red4] at (0,0.60) {$\alpha_0$};

        % Атомы m
        \foreach \x in {-5,-1,3}{
            \filldraw[
                fill=rc3,
                draw=rtext,
                line width=0.9pt
            ]
            (\x,0) circle (0.28);

            \node at (\x,0.60) {$m$};
        }

        % Атомы M
        \foreach \x in {-3,1,5}{
            \filldraw[
                fill=rbgsecondary,
                draw=rtext,
                line width=0.9pt
            ]
            (\x,0) circle (0.28);

            \node at (\x,0.60) {$M$};
        }

        % Индексы
        \node at (-5,-0.55) {$u_{-1}$};
        \node at (-3,-0.55) {$v_{-1}$};

        \node at (-1,-0.55) {$u_0$};
        \node at ( 1,-0.55) {$v_0$};

        \node at ( 3,-0.55) {$u_1$};
        \node at ( 5,-0.55) {$v_1$};

    \end{tikzpicture}

    \captionsetup{
        font=footnotesize,
        labelfont=bf,
        skip=5pt
    }

    \captionof{figure}{
        Двухатомная цепочка с дефектом связи.
        Розовая пружина между $u_0$ и $v_0$ имеет жёсткость
        $\alpha_0\neq\alpha$; остальные связи имеют жёсткость $\alpha$.
    }

    \label{fig:defect_spring_diatomic}
\end{rbox}


\solution

\subsection*{Уравнения движения в области дефекта}

Для центральной ячейки $n=0$
уравнение для атома $u_0$ имеет вид
%
\begin{equation*}
    -m\omega^2u_0
    =
    \alpha_0(v_0-u_0)
    -
    \alpha(u_0-v_{-1}).
\end{equation*}
%

Для атома $v_0$
%
\begin{equation*}
    -M\omega^2v_0
    =
    \alpha(u_1-v_0)
    -
    \alpha_0(v_0-u_0).
\end{equation*}
%

Для всех ячеек вдали от дефекта,
$n\neq0$,
сохраняются уравнения идеальной цепочки:
%
\begin{equation*}
    -m\omega^2u_n
    =
    \alpha(v_n-u_n)
    -
    \alpha(u_n-v_{n-1}),
\end{equation*}
%
\begin{equation*}
    -M\omega^2v_n
    =
    \alpha(u_{n+1}-v_n)
    -
    \alpha(v_n-u_n).
\end{equation*}


% ------------------------------------------------------------
% Локализованное решение
% ------------------------------------------------------------

\subsection*{Локализованное решение}

Ищем состояние,
амплитуда которого экспоненциально убывает
при удалении от дефекта:
%
\begin{equation*}
    u_n
    =
    u_0
    e^{-2\chi a|n|}
    (-1)^n,
\end{equation*}
%
\begin{equation*}
    v_n
    =
    v_0
    e^{-2\chi a|n|}
    (-1)^n.
\end{equation*}
%

Здесь
%
\begin{equation*}
    \chi>0
\end{equation*}
%
--- декремент пространственного затухания.

Множитель $(-1)^n$ описывает изменение фазы
на $\pi$ при переходе между соседними элементарными ячейками.

Для ближайших к дефекту атомов получаем
%
\begin{equation*}
    v_{-1}
    =
    -v_0e^{-2\chi a},
\end{equation*}
%
\begin{equation*}
    u_1
    =
    -u_0e^{-2\chi a}.
\end{equation*}


% ------------------------------------------------------------
% Центральная ячейка
% ------------------------------------------------------------

\subsection*{Система уравнений для центральной ячейки}

Подставляем найденные выражения в уравнения при $n=0$:
%
\begin{equation*}
    -m\omega^2u_0
    =
    \alpha_0(v_0-u_0)
    -
    \alpha
    \left(
        u_0+v_0e^{-2\chi a}
    \right),
\end{equation*}
%
\begin{equation*}
    -M\omega^2v_0
    =
    \alpha
    \left(
        -u_0e^{-2\chi a}-v_0
    \right)
    -
    \alpha_0(v_0-u_0).
\end{equation*}
%

Переносим все слагаемые в левую часть:
%
\begin{equation*}
    \begin{pmatrix}
        m\omega^2-(\alpha_0+\alpha)
        &
        \alpha_0-\alpha e^{-2\chi a}
        \\[2mm]
        \alpha_0-\alpha e^{-2\chi a}
        &
        M\omega^2-(\alpha_0+\alpha)
    \end{pmatrix}
    \begin{pmatrix}
        u_0
        \\
        v_0
    \end{pmatrix}
    =
    0.
\end{equation*}
%

Нетривиальное решение требует
%
\begin{equation*}
    mM\omega^4
    +
    (\alpha_0+\alpha)^2
    -
    \omega^2(m+M)(\alpha_0+\alpha)
    -
    \left(
        \alpha_0-\alpha e^{-2\chi a}
    \right)^2
    =
    0.
\end{equation*}


% ------------------------------------------------------------
% Уравнения вдали от дефекта
% ------------------------------------------------------------

\subsection*{Уравнения вдали от дефекта}

Вне центральной ячейки пружины имеют обычную жёсткость $\alpha$.

После подстановки затухающего решения
система в рукописном конспекте записана в форме
%
\begin{equation*}
    \begin{pmatrix}
        m\omega^2-2\alpha
        &
        \alpha(1-e^{-2\chi a})
        \\[2mm]
        \alpha(1-e^{2\chi a})
        &
        M\omega^2-2\alpha
    \end{pmatrix}
    \begin{pmatrix}
        u_n
        \\
        v_n
    \end{pmatrix}
    =
    0.
\end{equation*}
%

Для определения частоты нам нужен только определитель,
поэтому перестановка показателей
$e^{2\chi a}$ и $e^{-2\chi a}$
между внедиагональными элементами
на дальнейший результат не влияет.

Условие существования ненулевого решения:
%
\begin{equation*}
    \left(
        m\omega^2-2\alpha
    \right)
    \left(
        M\omega^2-2\alpha
    \right)
    -
    \alpha^2
    \left(
        1-e^{-2\chi a}
    \right)
    \left(
        1-e^{2\chi a}
    \right)
    =
    0.
\end{equation*}
%

Раскрываем скобки:
%
\begin{equation*}
    mM\omega^4
    -
    2\alpha(m+M)\omega^2
    +
    4\alpha^2
    -
    \alpha^2
    \left(
        2-e^{2\chi a}-e^{-2\chi a}
    \right)
    =
    0.
\end{equation*}
%

Используем
%
\begin{equation*}
    e^{2\chi a}+e^{-2\chi a}
    =
    2\cosh(2\chi a).
\end{equation*}
%

Тогда постоянная часть равна
%
\begin{align*}
    4\alpha^2
    -
    \alpha^2
    \left[
        2-2\cosh(2\chi a)
    \right]
    &=
    2\alpha^2
    \left[
        1+\cosh(2\chi a)
    \right]
    \\
    &=
    4\alpha^2
    \cosh^2(\chi a).
\end{align*}
%

Следовательно,
%
\begin{equation*}
    mM\omega^4
    -
    2\alpha(m+M)\omega^2
    +
    4\alpha^2
    \cosh^2(\chi a)
    =
    0.
\end{equation*}
%

После деления на $mM$ получаем
%
\begin{equation*}
    \omega^4
    -
    \frac{2\alpha(m+M)}{mM}
    \omega^2
    +
    \frac{4\alpha^2}{mM}
    \cosh^2(\chi a)
    =
    0.
\end{equation*}

% ------------------------------------------------------------
% Дискриминант
% ------------------------------------------------------------

\subsection*{Дискриминант}

Рассматриваем полученное выражение
как квадратное уравнение относительно
%
\begin{equation*}
    x=\omega^2.
\end{equation*}
%

Тогда
%
\begin{equation*}
    x^2
    -
    \frac{2\alpha(m+M)}{mM}x
    +
    \frac{4\alpha^2}{mM}
    \cosh^2(\chi a)
    =
    0,
\end{equation*}
%
а его дискриминант равен
%
\begin{equation*}
    D
    =
    \left[
        \frac{2\alpha(m+M)}{mM}
    \right]^2
    -
    4
    \frac{4\alpha^2}{mM}
    \cosh^2(\chi a).
\end{equation*}


\remark[Окончание записи.]
На этом месте присланная запись семинара обрывается.
Дальнейший вывод условия существования связанного состояния
и завершение доказательства его положения
в запрещённой зоне в исходных листах отсутствуют.


\end{document}
